\documentclass[10pt,letterpaper]{article}
\usepackage[margin=1in]{geometry}
\usepackage{listings}
\usepackage{xcolor}
\usepackage{titlesec}
\usepackage{parskip}
\usepackage{microtype}
\usepackage{enumitem}
\usepackage{multicol}
\usepackage{booktabs}
\usepackage[colorlinks=true,urlcolor=blue,linkcolor=black]{hyperref}

\definecolor{codebg}{RGB}{245,245,245}
\definecolor{codecomment}{RGB}{100,100,100}
\definecolor{codestring}{RGB}{180,50,50}
\definecolor{codekeyword}{RGB}{0,80,180}

\lstset{
  language=Python,
  backgroundcolor=\color{codebg},
  basicstyle=\ttfamily\small,
  keywordstyle=\color{codekeyword}\bfseries,
  stringstyle=\color{codestring},
  commentstyle=\color{codecomment}\itshape,
  showstringspaces=false,
  frame=single,
  framerule=0pt,
  rulecolor=\color{codebg},
  xleftmargin=6pt,
  xrightmargin=6pt,
  aboveskip=12pt,
  belowskip=4pt,
  breaklines=true,
  breakatwhitespace=true,
  tabsize=4,
  morekeywords={as, with, yield, lambda, None, True, False, self},
}

\titleformat{\section}{\large\bfseries}{}{0em}{}[\titlerule]
\titleformat{\subsection}{\normalsize\bfseries\itshape}{}{0em}{}
\titlespacing{\section}{0pt}{12pt}{4pt}
\titlespacing{\subsection}{0pt}{8pt}{2pt}

\begin{document}

\begin{center}
  {\LARGE \textbf{Python Idioms Reference}}\\[4pt]
  {\normalsize CodeSignal Assessment Prep}
\end{center}

\vspace{6pt}

%-----------------------------------------------------------------------
\section*{Quick Reference --- Open These Tabs First}

\begin{itemize}[noitemsep,leftmargin=*]
  \item \href{https://docs.python.org/3/library/functions.html}{\texttt{docs.python.org/3/library/functions.html}} --- Built-ins: \texttt{len, sorted, min, max, sum, enumerate, zip, range, isinstance, any, all, next}
  \item \href{https://docs.python.org/3/library/stdtypes.html#text-sequence-type-str}{\texttt{docs.python.org/3/library/stdtypes.html\#text-sequence-type-str}} --- String methods: \texttt{split, join, strip, find, replace, startswith}
  \item \href{https://docs.python.org/3/library/stdtypes.html#mapping-types-dict}{\texttt{docs.python.org/3/library/stdtypes.html\#mapping-types-dict}} --- Dict methods: \texttt{get, items, keys, values, setdefault, pop}
  \item \href{https://docs.python.org/3/library/threading.html}{\texttt{docs.python.org/3/library/threading.html}} --- \texttt{Thread, Lock, start, join}
  \item \href{https://docs.python.org/3/library/queue.html}{\texttt{docs.python.org/3/library/queue.html}} --- \texttt{Queue, put, get, Empty, task\_done}
  \item \href{https://docs.python.org/3/library/asyncio.html}{\texttt{docs.python.org/3/library/asyncio.html}} --- \texttt{async def, await, gather, run, sleep}
  \item \href{https://docs.python.org/3/library/collections.html}{\texttt{docs.python.org/3/library/collections.html}} --- \texttt{defaultdict, Counter, deque}
  \item \href{https://docs.python.org/3/library/unittest.html}{\texttt{docs.python.org/3/library/unittest.html}} --- \texttt{assertEqual, assertRaises, setUp, assertIsNone, assertTrue}
  \item \href{https://docs.python.org/3/library/stdtypes.html#list}{\texttt{docs.python.org/3/library/stdtypes.html\#list}} --- List methods: \texttt{append, extend, pop, remove, insert, sort}
\end{itemize}

\vspace{4pt}

%-----------------------------------------------------------------------
\section{Dict Patterns \normalfont\small\href{https://docs.python.org/3/library/stdtypes.html#mapping-types-dict}{docs}}

\begin{lstlisting}
d.get(key, default)           # safe lookup, no KeyError
d.get(key, 0) + 1             # increment with default -- memorize this
d.setdefault(key, []).append(x)  # get or create, then use

# Running counter (word count pattern)
counts = {}
for x in items:
    counts[x] = counts.get(x, 0) + 1

# Existence check
if key in d:        # True if key exists
if key not in d:    # True if key is absent -- prefer over "not key in d"
\end{lstlisting}

%-----------------------------------------------------------------------
\section{Sorting with Keys \normalfont\small\href{https://docs.python.org/3/library/functions.html#sorted}{docs}}

\begin{lstlisting}
sorted(items, key=lambda x: x["score"])                    # sort by field
sorted(items, key=lambda x: (-x["score"], x["name"]))      # desc score, asc name (tie-break)
sorted(items, key=lambda x: x[1], reverse=True)            # sort list of tuples by second element
lst.sort()        # sorts in place, returns None
sorted(lst)       # returns new list, original unchanged

# min/max without sorting the whole list
min(items, key=lambda x: x["id"])      # item with lowest id -- O(n), cleaner than sorted()[0]
max(items, key=lambda x: x["score"])   # item with highest score
\end{lstlisting}

%-----------------------------------------------------------------------
\section{Unpacking}

\begin{lstlisting}
a, b = b, a                   # swap without temp variable
first, *rest = [1, 2, 3, 4]   # first=1, rest=[2,3,4]
*init, last  = [1, 2, 3, 4]   # init=[1,2,3], last=4
a, b = some_tuple             # unpack a 2-tuple
for i, x in enumerate(lst):  # index + value simultaneously
for a, b in zip(list1, list2):  # parallel iteration

# Tuple unpacking in loops (common in CodeSignal action lists)
for action_type, *args in actions:    # action_type = first element, rest in args
    if action_type == "post":
        author, content = args
\end{lstlisting}

%-----------------------------------------------------------------------
\section{Comprehensions}

\begin{lstlisting}
[f(x) for x in lst if condition(x)]         # filter + transform
{k: v for k, v in d.items() if v > 0}       # filter a dict
{k: [] for k in keys}                        # initialize dict of lists
[[0]*cols for _ in range(rows)]              # 2D list (NOT [[0]*cols]*rows -- shared ref bug)

# Flattening one level
[x for sublist in nested for x in sublist]

# Set comprehension -- get unique values
{item["author"] for item in books}
\end{lstlisting}

%-----------------------------------------------------------------------
\section{String Operations \normalfont\small\href{https://docs.python.org/3/library/stdtypes.html#text-sequence-type-str}{docs}}

\begin{lstlisting}
", ".join(["a", "b", "c"])    # -> "a, b, c"
"hello world".split()         # -> ["hello", "world"] (any whitespace)
"a:b:c".split(":")            # -> ["a", "b", "c"] (specific delimiter)
"  hello  ".strip()           # -> "hello" (strips both ends)
"  hello  ".lstrip()          # -> "hello  " (left only)

# f-string format specs
f"{value:.2f}"                # float with 2 decimal places:  "3.14"
f"{value:>10}"                # right-align in 10-char field: "      3.14"
f"{value:<12}"                # left-align in 12-char field:  "3.14        "
f"{value:03d}"                # zero-pad integer:             "007"
f"{value!r}"                  # repr (adds quotes) -- useful for debugging

# Parsing a structured line (common in assessments)
fn_id, event, timestamp = log_line.split(":")
name, score, category = line.strip().split(":")

# Skip blank lines
lines = [l for l in text.split("\n") if l.strip()]
\end{lstlisting}

%-----------------------------------------------------------------------
\section{Control Flow Idioms}

\begin{lstlisting}
x = value if condition else default            # ternary
any(x > 0 for x in lst)                       # True if any element satisfies
all(x > 0 for x in lst)                       # True if all elements satisfy
next((x for x in lst if x > 0), None)         # first match or None (safe)

# Truthiness -- prefer over len() checks
if lst:           # True if list is non-empty
if not lst:       # True if list is empty
if d:             # True if dict is non-empty
if queued_jobs:   # cleaner than if len(queued_jobs) > 0:
\end{lstlisting}

%-----------------------------------------------------------------------
\section{Working with Indices}

\begin{lstlisting}
lst[-1]         # last element
lst[-2]         # second to last
lst[::-1]       # reversed copy
lst[i:j]        # slice i to j (exclusive)
(n-1)//2        # leftmost middle index (use for median -- verified against examples)
n//2            # rightmost middle index for even n
\end{lstlisting}

%-----------------------------------------------------------------------
\section{Collections (standard library) \normalfont\small\href{https://docs.python.org/3/library/collections.html}{docs}}

\begin{lstlisting}
from collections import defaultdict
d = defaultdict(int)       # missing keys return 0
d = defaultdict(list)      # missing keys return []
d["new_key"].append(x)     # no KeyError even if key is absent

from collections import Counter
c = Counter([1, 1, 2, 3])  # -> {1:2, 2:1, 3:1}
c.most_common(3)           # top 3 by count as list of (value, count) tuples

from collections import deque
q = deque()
q.appendleft(x)            # O(1) prepend (list.insert(0,x) is O(n))
q.popleft()                # O(1) pop from front
\end{lstlisting}

%-----------------------------------------------------------------------
\section{Class Design Patterns}

\subsection{Manager class -- standard CodeSignal pattern}
\begin{lstlisting}
class ManagerImpl:
    def __init__(self):
        self.items = {}       # id -> item dict -- O(1) lookup, not a list
        self.next_id = 1
        self.lock = threading.Lock()   # add lock immediately even if not yet needed

    def add_item(self, name):
        item_id = self.next_id
        self.items[item_id] = {
            "id": item_id, "name": name, "status": "pending",
            "total_ops": 0         # add ALL fields you'll need for later levels
        }
        self.next_id += 1
        return item_id

    def get_item(self, item_id):
        return self.items.get(item_id)   # returns None if not found -- use this
\end{lstlisting}

\subsection{Status transition pattern}
\begin{lstlisting}
def start(self, item_id):
    if item_id not in self.items:
        raise ValueError("not found")
    if self.items[item_id]["status"] != "pending":
        raise ValueError(f"expected pending, got {self.items[item_id]['status']}")
    self.items[item_id]["status"] = "active"
\end{lstlisting}

\subsection{Aggregation -- iterating stored items}
\begin{lstlisting}
# Total copies across all books
total = sum(b["copies"] for b in self.library.values())

# Unique users (flattening a list-of-lists field)
all_users = {u for item in self.items.values() for u in item["checked_out_by"]}

# Top n by field, with tie-breaking
top = sorted(self.items.values(),
             key=lambda x: (-x["total_checkouts"], x["title"]))[:n]
result = [{"id": b["id"], "title": b["title"],
           "total_checkouts": b["total_checkouts"]} for b in top]
\end{lstlisting}

%-----------------------------------------------------------------------
\section{Threading \normalfont\small\href{https://docs.python.org/3/library/threading.html}{docs}}

\begin{lstlisting}
import threading

# Basic thread -- start all, then join all (never join inside start loop)
threads = []
for f in funcs:
    t = threading.Thread(target=f)          # target is function object, no ()
    t.start()
    threads.append(t)
for t in threads:
    t.join()                                # waits for all to finish

# Thread with arguments
t = threading.Thread(target=worker, args=("Alice",))   # args must be a tuple
t = threading.Thread(target=worker, kwargs={"name": "Alice"})

# Lock -- prevents race conditions on shared state
lock = threading.Lock()
with lock:                                  # only one thread runs this block at a time
    shared_value += 1

# CAUTION: never hold a lock while calling t.join() on threads that need that lock
# That is a deadlock.
\end{lstlisting}

\subsection{Worker pool -- competing for jobs from shared state}
\begin{lstlisting}
# Pattern: lock wraps state ACCESS only -- work happens outside the lock
def run_parallel(self, n_workers):
    def worker():
        while True:
            with self.lock:
                job = self.start_next()    # atomically grab AND mark the job
            if job is None:
                break
            self.finish(job["id"])         # work outside lock = true parallelism

    threads = [threading.Thread(target=worker) for _ in range(n_workers)]
    for t in threads: t.start()
    for t in threads: t.join()

# KEY: start_next() must atomically find AND change status.
# If you separate "find" and "mark" into two locked calls, two workers can
# both see the same job between the two calls -- race condition.
\end{lstlisting}

\subsection{Ordered results from parallel tasks}
\begin{lstlisting}
# Pre-allocate results list -- write by index, not append
def parallel_map(tasks, n_workers):
    results = [None] * len(tasks)    # pre-allocate by length
    lock = threading.Lock()
    next_i = 0

    def worker():
        nonlocal next_i              # required to assign to enclosing variable
        while True:
            with lock:
                i = next_i
                if i >= len(tasks):
                    break
                next_i += 1          # grab index under lock (atomic check+increment)
            results[i] = tasks[i]()  # write to pre-allocated slot -- no lock needed
                                     # (each worker writes to a different index)

    threads = [threading.Thread(target=worker) for _ in range(n_workers)]
    for t in threads: t.start()
    for t in threads: t.join()
    return results

# nonlocal: required when an inner function ASSIGNS to an outer variable.
# Reading an outer variable is fine without it; assigning triggers UnboundLocalError.
def outer():
    count = 0
    def inner():
        nonlocal count    # without this: UnboundLocalError on count += 1
        count += 1
    inner()
    return count          # 1
\end{lstlisting}

\subsection{Thread-safe counter}
\begin{lstlisting}
class ThreadSafeCounter:
    def __init__(self):
        self._count = 0
        self._lock = threading.Lock()

    def increment(self):
        with self._lock:
            self._count += 1

    @property
    def value(self):
        with self._lock:
            return self._count
\end{lstlisting}

%-----------------------------------------------------------------------
\section{Queue (thread-safe communication) \normalfont\small\href{https://docs.python.org/3/library/queue.html}{docs}}

\begin{lstlisting}
import queue

q = queue.Queue()           # thread-safe FIFO -- no lock needed for get/put
q.put(item)                 # add item
item = q.get(block=False)   # raises queue.Empty if empty -- use in try/except
q.task_done()               # signal that a gotten item has been processed
q.join()                    # block until all items have been task_done()

# Worker pool with queue.Queue -- cleaner than manual lock when tasks are independent
def process_with_workers(items, fn, n_workers):
    q = queue.Queue()
    results = {}
    lock = threading.Lock()    # queue is thread-safe; results dict is NOT

    for item in items:
        q.put(item)

    def worker():
        while True:
            try:
                item = q.get(block=False)    # raises queue.Empty when done
            except queue.Empty:
                break
            value = fn(item)
            with lock:
                results[item] = value        # protect the results dict

    threads = [threading.Thread(target=worker) for _ in range(n_workers)]
    for t in threads: t.start()
    for t in threads: t.join()
    return results

# Use queue.Queue when distributing plain work items.
# Use manual lock when "get next item" has side effects (e.g. marking a job as running).
\end{lstlisting}

%-----------------------------------------------------------------------
\section{asyncio \normalfont\small\href{https://docs.python.org/3/library/asyncio.html}{docs}}

\begin{lstlisting}
import asyncio

# Declare a coroutine
async def fetch(url):
    await asyncio.sleep(1)      # await suspends this coroutine, lets others run
    return url

# Run concurrently -- equivalent of threading run_parallel
async def main():
    results = await asyncio.gather(fetch("a"), fetch("b"), fetch("c"))
    return results

asyncio.run(main())             # entry point -- starts the event loop, call once

# Gather from a list (unpack with *)
coros = [fetch(url) for url in urls]
results = await asyncio.gather(*coros)

# asyncio.sleep vs time.sleep
await asyncio.sleep(1)          # yields control to event loop (correct)
time.sleep(1)                   # BLOCKS entire event loop (wrong in async code)
\end{lstlisting}

%-----------------------------------------------------------------------
\section{Threading vs asyncio}

\begin{center}
\begin{tabular}{lll}
\toprule
 & \textbf{threading} & \textbf{asyncio} \\
\midrule
Syntax & regular functions & \texttt{async def / await} \\
Parallelism & OS threads (true concurrency) & single thread, cooperative \\
Shared state & needs Lock & safer by default (one thing runs at a time) \\
Best for & any blocking work, simpler mental model & high-concurrency I/O \\
Worker loop & natural \texttt{while True} & needs \texttt{await} to yield \\
CodeSignal advice & \textbf{prefer this} & use if spec requires it \\
\bottomrule
\end{tabular}
\end{center}

%-----------------------------------------------------------------------
\section{Exception Handling \normalfont\small\href{https://docs.python.org/3/library/exceptions.html}{docs}}

\begin{lstlisting}
try:
    result = risky_operation()
except KeyError:
    result = default
except (TypeError, ValueError) as e:
    raise RuntimeError("bad input") from e
except Exception:               # catches all normal errors (not KeyboardInterrupt)
    return None
finally:
    cleanup()                   # always runs

# Raise with message
raise ValueError("Cannot divide by zero")
raise TypeError(f"Expected int, got {type(x).__name__}")

# isinstance for type checking
isinstance(x, (int, float))     # True if x is int or float -- not type(x) == int
\end{lstlisting}

%-----------------------------------------------------------------------
\section{In-Memory Key-Value Store Patterns}

\subsection{Basic structure with TTL (time-to-live) \normalfont\small\href{https://docs.python.org/3/library/time.html}{time docs}}
\begin{lstlisting}
import time

class KVStore:
    def __init__(self):
        self.data    = {}   # key -> value
        self.expires = {}   # key -> expiry timestamp (float)
                            # two separate dicts -- cleaner than nesting

    def set(self, key, value):
        self.data[key] = value
        self.expires.pop(key, None)   # clear any existing TTL

    def set_with_ttl(self, key, value, ttl):
        self.data[key] = value
        self.expires[key] = time.time() + ttl   # store absolute expiry time

    def _expired(self, key):
        if key not in self.expires:
            return False
        return time.time() >= self.expires[key]   # >= not >

    def get(self, key):
        if key not in self.data or self._expired(key):
            return None
        return self.data[key]

    def delete(self, key):
        if key in self.data and not self._expired(key):
            del self.data[key]
            self.expires.pop(key, None)
            return True
        return False

# Prefix scan -- return all live keys matching a prefix
def scan(self, prefix):
    return {k: v for k, v in self.data.items()
            if k.startswith(prefix) and not self._expired(k)}
\end{lstlisting}

\subsection{Transactions (begin / commit / rollback)}
\begin{lstlisting}
_DELETED = object()   # sentinel object -- distinct from any real value

class KVStore:
    def __init__(self):
        self.data   = {}
        self.buffer = None   # None = no active transaction; dict = transaction open

    def begin(self):
        self.buffer = {}

    def commit(self):
        if self.buffer is None:
            return False
        for key, value in self.buffer.items():
            if value is _DELETED:        # buffered deletion
                self.data.pop(key, None)
            else:
                self.data[key] = value
        self.buffer = None
        return True

    def rollback(self):
        if self.buffer is None:
            return False
        self.buffer = None               # discard without applying
        return True

    def set(self, key, value):
        if self.buffer is not None:
            self.buffer[key] = value     # buffer it -- don't touch data yet
        else:
            self.data[key] = value

    def get(self, key):
        if self.buffer is not None and key in self.buffer:
            v = self.buffer[key]
            return None if v is _DELETED else v   # respect buffered deletes
        return self.data.get(key)

    def delete(self, key):
        if self.buffer is not None:
            if key in self.buffer or key in self.data:
                self.buffer[key] = _DELETED
                return True
            return False
        if key in self.data:
            del self.data[key]
            return True
        return False

# Pattern summary:
#   self.buffer = None   ->  no transaction
#   self.buffer = {}     ->  transaction open; writes go here
#   _DELETED sentinel    ->  records a deletion inside the buffer
#   get() checks buffer first, then data  ->  reads see uncommitted writes
\end{lstlisting}

%-----------------------------------------------------------------------
\section{Stack Simulation Pattern}

\begin{lstlisting}
# Pattern: prev_time tracks when current top-of-stack last started or resumed.
# Used for problems involving nested function calls, intervals, or nested events.
stack = []
prev_time = 0

for log in logs:
    fn_id, event, timestamp = log.split(":")
    fn_id, timestamp = int(fn_id), int(timestamp)

    if event == "start":
        if stack:
            exclusive[stack[-1]] += timestamp - prev_time   # time before this call
        stack.append(fn_id)
        prev_time = timestamp

    else:  # end
        exclusive[stack[-1]] += timestamp - prev_time + 1  # +1: end is inclusive
        stack.pop()
        prev_time = timestamp + 1                           # resume at next unit
\end{lstlisting}

%-----------------------------------------------------------------------
\section{R-to-Python Gotchas}

\begin{lstlisting}
# Indexing: Python is 0-indexed
lst[0]           # first element (not lst[1])
lst[-1]          # last element (not lst[length])

# Integer division
10 / 3           # -> 3.333...  (float, even for ints)
10 // 3          # -> 3         (floor division)

# None check
if x is None:    # correct
if x == None:    # works but wrong style

# Mutability: Python passes lists by reference
def f(lst):
    lst.append(1)  # modifies the ORIGINAL list

# append returns None -- don't assign it
result = lst.append(x)   # result is None -- common bug

# Logical operators
not x    # not !x
and      # not &&
or       # not ||

# No paste() -- use f-strings
f"hello {name}, you are {age} years old"

# range vs R's seq
range(n)          # 0, 1, ..., n-1  (not 1 to n)
range(a, b)       # a, a+1, ..., b-1

# Shadowing builtins -- don't use these as variable names
# id, list, dict, set, type, input, range, min, max, sum, len, zip
\end{lstlisting}

%-----------------------------------------------------------------------
\section{Assessment Strategy}

\begin{itemize}[noitemsep]
  \item Read the \textbf{full spec} before writing any code -- all levels.
  \item Open the test file first -- tests are the spec, not the description.
  \item \textbf{Sketch \texttt{\_\_init\_\_} before touching code}: what data structures, what fields does Level 4+ require? Add all fields upfront.
  \item Use \texttt{dict} keyed by id, not a list -- O(1) lookup, simpler methods.
  \item Start with Level 1, pass all tests before moving on.
  \item Run tests early and often (\texttt{bash run\_single\_test.sh "<name>"} for one test).
  \item Don't optimize. Write correct first.
  \item For concurrency: use \texttt{threading} with \texttt{Lock} unless \texttt{asyncio} is explicitly required.
  \item If stuck on a level, move on -- partial credit beats zero.
  \item Leave 5 minutes to re-run all tests at the end.
  \item Language docs and Google (AI Overview only) are allowed.
\end{itemize}

\end{document}
